\documentclass{article}
\usepackage[preprint]{neurips_2026}

\usepackage[utf8]{inputenc}
\usepackage[T1]{fontenc}
\usepackage{hyperref}
\usepackage{url}
\usepackage{booktabs}
\usepackage{amsfonts}
\usepackage{amsmath}
\usepackage{amssymb}
\usepackage{microtype}
\usepackage{xcolor}
\usepackage{graphicx}

% This is an ES 678 course paper, not a conference submission.
\makeatletter
\renewcommand{\@noticestring}{ES 678 Course Project}
\makeatother

\title{How Small Can a Trustworthy Prediction Set Be?\\ Score Design and Class-Conditional Coverage in Conformal Classification}

\author{%
  <Your Name> \\
  Roll Number: <Roll Number> \\
  ES 678 \\
  Indian Institute of Technology Gandhinagar
}

\begin{document}
\maketitle

\begin{abstract}
% 100--150 words: problem, data, method, evaluation, main finding.
% Draft below; update numbers once results exist.
We study how a fixed image classifier can produce trustworthy prediction \emph{sets} at a
controlled error rate using conformal prediction, and how the choice of conformity score affects
efficiency and conditional coverage. The marginal coverage guarantee is distribution-free and
holds for \emph{any} score, so the research question is necessarily about set size and coverage
beyond the marginal average. We compare naive, adaptive (APS) and regularised (RAPS) scores, a
temperature-scaled variant, a proposed Mahalanobis-distance-augmented score, and class-conditional
(Mondrian) conformal prediction, on MNIST and CIFAR-10, with an out-of-distribution failure audit.
We derive the coverage theorem and quantify the efficiency--conditional-coverage frontier.
\end{abstract}

\section{Problem and Motivation}
State the question in concrete terms: for a fixed classifier, how do score design and feature-space
distance affect set size at fixed coverage, and the class-conditional coverage that marginal
guarantees fail to control. Define the target (a label set $C(x)$), what one row represents (one
image $+$ label), and when a prediction is made (inference time, no label available). Note the
limits of the claim: distribution-free \emph{marginal} coverage under exchangeability; no guarantee
under distribution shift.

\section{Data and Preparation}
Sources, sizes, and access: MNIST, CIFAR-10 (in-distribution); FashionMNIST, KMNIST, SVHN,
CIFAR-100 (OOD); all via \texttt{torchvision}. Describe the four-way split
($D_{train}, D_{design}, D_{cal}, D_{test}$) and state explicitly that all fitted preprocessing and
all tuning use training/design data only.

\section{Methods}
Define the baseline (naive score, top-1) and each method: APS, RAPS, temperature scaling,
Mahalanobis distance, the proposed $s_\gamma(x,y) = (1 - p_T(y|x)) + \gamma\,\tilde d_M(x,y)$, and
Mondrian conformal. Give the objective/assumptions and define every symbol.

\section{Evaluation Design}
Describe the split before reporting results; separate tuning from final testing; define each metric
(marginal coverage, mean/median/90th-percentile set size, singleton/empty rates, class-conditional
coverage, worst-class gap, OOD behaviour). Justify each metric.

\section{Results and Error Analysis}
Main table (coverage vs size per score, per $\alpha$) and the efficiency frontist figure. Interpret
in context; inspect failure cases (which classes/inputs break coverage); report uncertainty across
seeds and subgroup results.

\section{Graduate Depth Analysis}
Present the coverage theorem derivation with the two-sided bound
$1-\alpha \le P(Y \in C(X)) \le 1-\alpha + 1/(n+1)$, explain why marginal coverage cannot imply
conditional coverage, and report the conditional-coverage audit, the $\alpha \times \gamma$
sensitivity study, and the OOD failure-mode audit.

\section{Limitations and Conclusion}
Most important limitation and what evidence would strengthen the claim. Do not interpret predictive
association as proof of causation. Conclude with the main answer.

\section{Reproducibility and AI-Use Disclosure}
Code and data links/instructions, split details, seeds. State which AI tools were used, how, and
which outputs were independently verified.

\bibliographystyle{plainnat}
\bibliography{refs}

\end{document}
